% Vietnamese translation for OI Public Library
% Source-File: IOI中国国家候选队论文/2005/王俊/王俊.ppt
\documentclass[11pt]{article}
\usepackage{oipl-vn}

\newcommand{\Oh}{\mathcal{O}}

\title{Bản trình chiếu: Ghép cặp đồ thị hai phía}
\author{Wang Jun}
\date{2005}

\begin{document}
\maketitle

\origtitle{Applications of bipartite matching}
\sourcepath{IOI China candidate team papers / 2005 / Wang Jun / Wang Jun.ppt}

\section{Phạm vi bản dịch}

Bản trình chiếu đi cùng bài viết về ứng dụng ghép cặp hai phía. Phần chữ đọc
được từ PPT có mốc \textit{Roads}, các ký hiệu
\[
  n=|V_0|,\qquad m=|E_0|,
\]
và nhiều hình đồ thị không có lớp văn bản. Bản ghi chú này dùng bài viết đã
dịch làm neo, khôi phục mạch slide: nhận ra quan hệ một--một, dựng đồ thị hai
phía, rồi chọn thuật toán ghép phù hợp.

\section{Từ quan hệ một--một đến matching}

Ghép cặp hai phía xuất hiện khi quyết định của bài toán có dạng chọn cặp giữa
hai loại đối tượng. Mỗi đối tượng ở một phía chỉ được dùng nhiều nhất một lần.
Ta dựng một đồ thị hai phía \(X\cup Y\); cạnh \((x,y)\) nghĩa là \(x\) có thể
ghép với \(y\). Khi đó:
\begin{itemize}
  \item tìm được nhiều cặp nhất là ghép cực đại;
  \item trong số các ghép lớn nhất, tối ưu tổng trọng số là ghép trọng số;
  \item kiểm tra tồn tại một phương án đầy đủ là kiểm tra matching hoàn hảo.
\end{itemize}

Điều quan trọng nhất của slide là bước mô hình hóa. Thuật toán đường tăng hay
KM chỉ phát huy tác dụng sau khi hai phía và cạnh đã được chọn đúng.

\section{\textit{Railway Communication}}

Ví dụ đầu trong bài viết là \textit{Railway Communication}. Sau khi phân tích,
bài toán trở thành chọn các đoạn/đường sao cho mỗi đối tượng chỉ nhận một
nhiệm vụ. Cấu trúc này khớp trực tiếp với matching: một phía biểu diễn các
nhiệm vụ hoặc đoạn cần phục vụ, phía kia biểu diễn tài nguyên có thể dùng, và
một cạnh biểu diễn khả năng gán hợp lệ.

Nếu mục tiêu chỉ là phủ được nhiều nhất, ta dùng ghép cực đại. Nếu còn có chi
phí, ta dùng ghép trọng số. Bài viết chuyển bài toán trọng số nhỏ nhất sang
dạng trọng số lớn nhất bằng cách đổi trọng số:
\[
  W'=C-w,
\]
với \(C\) đủ lớn để bảo toàn ưu tiên về kích thước ghép. Nhờ đó có thể dùng
KM cho ghép trọng số lớn nhất.

\section{\textit{Roads}}

Mốc \textit{Roads} trong PPT là ví dụ quan trọng hơn. Bài toán gốc liên quan
đến đồ thị \(G_0=(V_0,E_0)\), với \(n=|V_0|\) và \(m=|E_0|\). Mục tiêu có thể
nhìn như một bài inverse MST: điều chỉnh hoặc chọn thông tin sao cho một cây
khung mong muốn trở nên tối ưu.

Bài viết đưa bài toán về ghép cặp trọng số. Một phía của đồ thị hai phía biểu
diễn các lựa chọn cần được gán, phía kia biểu diễn các cạnh/điều kiện có thể
đảm nhận lựa chọn ấy; trọng số cạnh ghi mức lợi ích hoặc mức tiết kiệm. Khi
chạy KM, tổng trọng số ghép lớn nhất tương ứng với phương án tối ưu trong bài
gốc.

Nếu chỉ cần giá trị mà không cần khôi phục phương án theo đúng dạng matching,
bài viết cũng chỉ ra cách nhìn như dòng cực đại chi phí lớn nhất: thêm nguồn,
đích, các cung dung lượng \(1\), rồi tìm đường tăng có chi phí lớn nhất trong
mạng dư. Đây là cùng một bản chất với matching trọng số, nhưng đôi khi tiện
hơn để tối ưu cài đặt.

\section{Tối ưu thuật toán}

Slide có nhiều sơ đồ hình ảnh, nhiều khả năng minh họa đường tăng. Về mặt kỹ
thuật, khi đồ thị lớn, không nên dựng và quét toàn bộ cạnh một cách máy móc.
Ta cần khai thác cấu trúc của bài:
\begin{itemize}
  \item loại bỏ các cạnh ảo chỉ dùng để ép matching hoàn hảo nếu chúng không
  ảnh hưởng đến giá trị;
  \item chỉ duy trì phần đồ thị còn có thể tham gia đường tăng;
  \item chuyển sang mô hình dòng chi phí nếu mạng dư dễ cập nhật hơn;
  \item dùng tính chất của \(G_0\) để sinh cạnh hai phía theo nhu cầu.
\end{itemize}

\section{Kết luận}

Ứng dụng matching trong thi tin học không chỉ là nhận ra tên thuật toán. Cần
phân tích ràng buộc một--một, dựng hai phía đúng, gán trọng số đúng, rồi chọn
giữa ghép cực đại, KM hoặc dòng chi phí. Hai ví dụ của slide cho thấy cùng một
bài toán có thể được nhìn theo nhiều mô hình tương đương; mô hình tốt là mô
hình làm phần cài đặt và chứng minh tối ưu ngắn nhất.

\end{document}
